\documentclass[journal, a4paper]{IEEEtran}
\usepackage{graphicx} 
\usepackage{tabularx}
\usepackage{tabulary}
\usepackage{url}       
\usepackage{amsmath} 
\usepackage[left=2cm,right=1.5cm,top=1.5cm,bottom=1.5cm]{geometry}
\usepackage{amsfonts}
\usepackage{amssymb}
\usepackage{textcomp,gensymb}
\usepackage{gensymb}
\usepackage{float}
\parindent=5mm 
\renewcommand{\tablename}{Tabla}
\renewcommand{\figurename}{Figura}
\renewcommand{\refname}{Referencias}

\begin{document}

% Define  el titulo del documento y el autor
	\title{Transferencia de calor en flujos laminares}
	\author{Acosta, Jimena; Arguello, Melissa; De Baca, Juan Carlos; Fleitas, Jazmin; Rojas, Mariana; Romero, Rosa; Strubing Yazmin; Vera, Gast\'on 
    \thanks{}}
	\markboth{LABORATORIO DE INGENIERIA DE ALIMENTOS II - FACULTAD DE CIENCIAS QUIMICAS - UNA}{}
	\maketitle

% Write abstract here
\begin{abstract}
Se determinan las variables relacionadas con las operaciones de transferencia de calor en flujos laminares de forma experimental para un sistema agua/aceite en r\'egimen estacionario. Estas variables son la tasa de transferencia de calor, diferencia logar\'itmica de temperatura media, el coeficiente global de transferencia de calor, los coeficientes de pel\'icula dentro y fuera de los tubos, viendo el efecto de la velocidad del fluido en estas variables. Adem\'as se compara el perfil de temperatura de la disposici\'on en contra-corriente, con la disposici\'on en co-corriente.  Al final se busca establecer la relaci\'on entre los n\'umeros de Nusselt y Graetz en tuber\'ias de secci\'on circular.
\end{abstract}

% Each section begins with a \section{title} command
\section{Introducci\'on}
Para calcular la cantidad de energ\'ia perdida debido a la fricci\'on en un sistema de fluido, es necesario caracterizar la naturaleza del flujo. Un flujo lento y uniforme se conoce como flujo laminar, mientras que un flujo r\'apido y ca\'otico se conoce como flujo turbulento. Los m\'etodos que se utilizan para calcular la p\'erdida de energ\'ia es diferente para cada tipo de flujo \cite{Mott}.

En el flujo laminar, s\'olo hay transferencia de calor por conducci\'on, puesto que no existen remolinos que lleven calor a la superficie isot\'ermica. Cuando el fluido se acerca a la superficie de calentamiento, es posible que tenga ya una capa l\'imite hidrodin\'amica total o parcialmente desarrollada. O quiz\'a se acerque a la superficie de calentamiento a una velocidad uniforme y las dos capas comiencen a formarse al mismo tiempo. La velocidad del fluido a trav\'es del tubo
y en todos los puntos de una secci\'on transversal cualquiera de la corriente es constante, la temperatura de la pared es constante, y las propiedades del fluido son independientes de la temperatura \cite{McCabe}.

Existen dos tipos de arreglos para un intercambaidor de calor en cuanto a su flujo, que corresponden a un arreglo en co-corriente y contracorriente. De la misma forma, se puede trabajar con flujo paralelo o cruzado, esto corresponde al \'angulo de ataque del fluido al tubo. Se obtiene un mejor rendimiento en cuanto a la transferencia de calor para un flujo en contracorriente, mientras que la virtud de un sistema en co-corriente reside en la mantenci\'on de una temperatura relativamente constante entre los fluidos. 

% Main Part
\section{Materiales y metodolog\'ia}
 El equipo utilizado para la pr\'actica es el intercambiador de calor de flujo laminar agua/aceite, es b\'asicamente un intercambiador de calor de tubos conc\'entricos con aceite caliente fluyendo por el tubo central y el agua de enfriamiento que fluye por el anular.
    Un termopar digital tipo selector muestra la temperatura de aceite caliente, agua fr\'ia y de la pared de metal registrada por los distintos termopares. Los circuitos de aceite y de agua fr\'ia son como sigue:
    
    
Circuito de aceite caliente.


El aceite caliente de un tanque de calentamiento equipado con un calentador de tipo resistencia el\'ectrica es bombeado por una bomba centr\'ifuga a trav\'es de una tuber\'ia al extremo superior del intercambiador en el tubo central. 
El aceite es enfriado conforme fluye hacia abajo a trav\'es del intercambiador, al salir pasa a trav\'es de una v\'alvula de control a un tanque de vidrio que posee una escala que permite la medici\'on del contenido y est\'a equipado con una v\'alvula de acci\'on r\'apida que permite al aceite volver al tanque de calentamiento.


Circuito de agua fr\'ia.


La tuber\'ia de agua fr\'ia pasa a trav\'es de una v\'alvula de control de flujo y un fluj\'ometro, luego va a una uni\'on de aflojamiento r\'apido. Una uni\'on similar de aflojamiento r\'apido conecta los extremos con el punto de drenaje. Las conexiones son totalmente reversibles de manera que el agua de enfriamiento puede fluir en direcci\'on co-corriente evitando que las mangueras se crucen o en contra-corriente haciendo cruzar las mangueras. 

El control de la temperatura del aceite se hace mediante una perilla de control del calentador que se encuentra en el centro del panel. Una luz piloto de ne\'on encima de la perilla destellar\'a en proporci\'on al ingreso de calor, cuando la temperatura de aceite a la entrada supera los $89^{\circ}$C el termostato apagar\'a el suministro el\'ectrico autom\'aticamente \cite{edibon}.

\setlength{\parskip}{0.1 pt}

\section{Resultados y discusi\'on}
	 En cuanto a la determinaci\'on de la tasa de transferencia de calor tanto para el agua como para el aceite, se observa un aumento lineal con el aumento del flujo del aceite, con un ajuste de $r=0.91$, lo mismo ocurre con el coeficiente general de calor, con un ajuste $r=0.99$.
En la tabla 1, se observa el aumento del coeficiente de transferencia de calor de superficie dentro y fuera del tubo al aumentar las temperaturas.
Al comparar la transferencia de calor en con flujo en contra corriente, se observa una eficaz transferencia de calor a medida que el flujo de agua avanza hacia la salida. En cuanto al flujo en co-corriente, la transferencia de calor no es la esperada, esto pudo deberse a la elevada temperatura a la cual estaba operando el equipo, lo que dificulto la lectura correcta de los termopares.
La relaci\'on entre el numero de Nusselt y Graetz posee un comportamiento similar a la formula correlativa, aunque los valores no son aproximados, como se observa en la figura 1.

\setlength{\intextsep}{0.3 pt} %distancia entre los flotadores insertados dentro del texto de la página (utilizando h) y el texto propiamente dicho.
 
	\begin{table}[H] %tabla 1
		\begin{center}
        \resizebox*{8,5 cm}{!}{
		\begin{tabular}{rcccccc}
			\hline
            \textbf Qagua (kW) & Qaceite (kW) & v (m/s) & U (W/m$^2$) & $\Delta$Tml aceite ($^{\circ}$C) & $\Delta$Tml agua ($^{\circ}$C)		\\	\hline
0,0627  & 0,0630   & 0,0263   & 48,2050  & 33,0151  & 3,2844 \\ \hline
0,0878  & 0,0748   & 0,0581   & 50,9077  & 37,6230  & 3,7574 \\ \hline
0,1128  & 0,0872   & 0,1013   & 57,2363  & 39,2339  & 3,9959 \\ \hline
0,1881  & 0,1403   & 0,3399   & 91,2507  & 39,2775  & 4,5896  \\ \hline
0,2006  & 0,1944   & 0,6022   & 128,2213 & 37,3970  & 4,8775  \\ \hline
0,1881  & 0,1929   & 0,6754   & 129,9099  & 36,2468 & 4,8775  \\ \hline
0,1881  & 0,1948   & 0,6847   & 134,3197  & 35,2439 & 4,8775  \\ \hline
0,1254  & 0,1293   & 0,2687   & 92,2011  & 34,5505   & 4,2843 \\ \hline			\end{tabular}}
		\caption{Determinaci\'on de la tasa de transferencia de calor, diferencia de temperatura media logaritmica y coeficiente global de transferencia de calor}
		\label{tabla:parte1}
		\end{center}
	\end{table}
    

   \begin{table}[!ht] %tabla 2
   		\begin{center}
        \resizebox*{8,5 cm}{!}{
		\begin{tabular}{rccccc}
		\hline
      	\textbf Tml aceite ($^{\circ}$C)  & $\Delta$Tml agua ($^{\circ}$C) & haceite (W/m$^2{\circ}$ C) & hagua (W/m $^2{\circ}$C) \\ \hline

33,0151 & 3,2843 & 53,0252  & 530,2941  & 48,2051    \\ \hline
37,6230 & 3,7574 & 55,2413  & 648,9396  & 50,9077    \\ \hline
39,2339 & 3,9959 & 61,7406  & 784,5448  & 57,2363    \\ \hline
39,2775 & 4,5896 & 99,2022  & 1138,4388 & 91,2507    \\ \hline
37,3970 & 4,8775 & 144,4280 & 1142,6534 & 128,2213   \\ \hline
36,2468 & 4,8775 & 147,8384 & 1071,2375 & 129,9099   \\ \hline
35,2439 & 4,8775 & 153,5762 & 1071,2375 & 134,3197   \\ \hline
34,5505 & 4,2843 & 103,9943 & 813,0387  & 92,2011    \\ \hline
		\end{tabular}}
        \caption{Determinaci\'on del coeficiente de transferencia de calor de superficie dentro y fuera del tubo y el efecto de la velocidad del fluido en estos}
        \label{tabla:parte2}
        \end{center}
	\end{table}

	\begin{figure}[!ht]
		\begin{center}
		\includegraphics[width=\columnwidth]{1.jpg}
		\caption{Efecto de la velocidad de flujo en el tubo interno con el coeficiente global de transferencia de calor y los coeficientes de pel\'icula}
		\label{fig:1}
		\end{center}
	\end{figure}
    

\begin{table}[!ht ]
\begin{center}
\begin{tabularx}{8,5 cm}{p{4,25 cm} p{4,25 cm}}
\hline
\textbf {Temperaturas ($^{\circ}$C)}	 \\ \hline
Contracorriente      & Co-corriente     \\ \hline
70,6                 & 72,4             \\ \hline
57,1                 & 70,8             \\ \hline
26,1                 & 29,7             \\ \hline
33,4                 & 33,4             \\ \hline
25,3                 & 24,6             \\ \hline
24,8                 & 26,9             \\ \hline
\end{tabularx}
\caption{Comparaci\'on del flujo concurrente y contra-corriente}
\label{tabla:parte3}
\end{center}
\end{table}

\begin{figure}[H]
	\begin{center}
		\includegraphics[width=\columnwidth]{2.jpg} 
	\end{center}
	\caption{Distribuci\'on de temperatura en contracorriente}
\end{figure}

\begin{figure}[H]
	\begin{center}
		\includegraphics[width=\columnwidth]{3.jpg} 
	\end{center}
	\caption{Distribuci\'on de temperatura en co-corriente}
\end{figure}

\begin{table}[H] %tabla 3
   	\begin{center}
		\begin{tabularx}{8,5 cm}{p{2,83 cm} p{2,83 cm} l}
		\hline
        \textbf Nu & Gz & Nu promedio \\ \hline
        0,00431068 	& 3,6907.10$^{-5}$ 	& 0.0582494 	\\ \hline
        0.004490838 & 8,2154.10$^{-5}$ 	& 0.07607614 	\\ \hline
        0,005019197 & 0,00014363 		& 0,09164678	\\ \hline
		0,008064641 & 0,00048215		&  0,13722484 	\\ \hline
		0,011741274 &  0,00085401		&  0,16603214	 \\ \hline
		0,012018522 & 0,0009532  		&  0,17222627  \\ \hline
 		0,012484972 & 0,0009628		  & 0,17280231		 \\ \hline
		0,008454216 & 0,00037874 		&  0,12661551	\\ \hline 
        \end{tabularx}
        \caption{Establecimiento de la relaci\'on entre n\'umeros de Nussel y Graetz para flujo laminar}
      \label{tab:parte3}
        \end{center}
        \end{table}
        
 \begin{figure}[!hb]
	\begin{center}
		\includegraphics[width=\columnwidth]{4.jpg} 
	\end{center}
	\caption{Comparaci\'on del n\'umero de Nussel promedio y n\'umero de Nussel calculado mediante correlaci\'on con relaci\'ion al n\'umero de Graetz}
\end{figure}
            
\section{Conclusi\'on}
	El aumento de temperatura del fluido fr\'io obtenido para los distintos caudales de fluido caliente ingresados, no fueron significativos, lo que comprueba que la transferencia de calor de un flujo laminar es ineficiente.

La tasa de transferencia de calor del fluido caliente y del fluido fr\'io aumentan linealmente con el incremento de la tasa de flujo m\'asico del aceite haciendo m\'as efectivo el intercambio de calor debido al aumento de la velocidad de entrada del fluido caliente. 

El coeficiente de pel\'icula correspondiente al flujo caliente es el que tiene mayor incidencia en el coeficiente global de transferencia de calor, el aumento de flujo del fluido caliente aumenta el valor del coeficiente global de transferencia de calor.

La transferencia de calor es m\'as efectiva cuando la configuraci\'on de flujo es en contra-corriente en lugar de co-corriente.


\begin{thebibliography}{5}

		\bibitem{Mott} % libro
        Mott R.L., Mec\'anica de Fluidos Aplicada. Prentise Hall Hispanoamericana, M\'exico, 4ta ed. 1996.
        \bibitem{McCabe} % libro
		W. McCabe, J.Smith, P. Harriot – Operaciones unitarias en ingenier\'ia qu\'imica. McGraw-Hill Interamericana, M\'exico, 7ma ed. 2007.
		\bibitem{edibon} %libro
    	Edibon. Manual experimental de operaci\'on y mantenimiento. España, 1998.

\end{thebibliography}

\end{document}